\subsection{Narrowing Down an Estimate} 
In general, we cannot know the \makeblank[1in] value of an irrational root. But we can get as close as we like.
\begin{itemize}
\item How do we describe how close we want to be?
\item How do we know when we got there?
\end{itemize}
\begin{Example}
The polynomial $x^5+x-42$ has a single real zero between $2$ and $3$. Find an approximation of this zero to within $0.01$.
\begin{itemize}
\item The fastest way to narrow the range is to cut it in \makeblank[1in] at each step.
\item We keep narrowing until we have points that are \makeblank apart.
\item Then our estimate is \makeblank them.
\end{itemize}
\begin{icpic}
\draw[|<->|,thick] (-1,0) node[anchor=north]{$a$} -- (1,0) node[anchor=north]{$b$};
\end{icpic}
We'll use our calculator to help us.
\begin{cpic}{}{8}
\foreach \x/\lbl in {-8/{$a$ (left)},-4/{$b$ (right)},0/Distance,4/{$c$ (midpoint)},8/$f(c)$}
	\regtext{\x}{-1}{\lbl};
\end{cpic}
The number \makeblanksmall is within 0.01 of the zero.
\begin{itemize}
\item We might write $x=\makeblank[1in]$
\item If we want the number to look rounded, we might get \makeblanksmall or \makeblanksmall.
\begin{itemize}
\item If the zero is between \makeblanksmall and \makeblanksmall, we would round it to \makeblanksmall.
\item If the zero is between \makeblanksmall and \makeblanksmall, we would round it to \makeblanksmall.
\item We just need to check $f(\makeblanksmall)=\makeblanksmall$.
\item We could also estimate the zero as \makeblanksmall.
\end{itemize}
But this extra step is unnecessary.
\end{itemize}
\end{Example}